\section{Introduction}

This paper gives a recursive obstruction theorem for projective
Reed--Solomon deep holes at arbitrary redundancy and exact classifications
through redundancy ten.  The common mechanism is a coherent polar flag, a
marked contraction that propagates splitting information without forgetting
the factor needed for a squarefree lift.  The contained part of the
recursion is exactly the persistent catalecticant scheme together with one
maximal Lucas carrier.  At a fixed depth, the transverse part escapes by a
uniform point count once the explicit one-step lower packages at every
intermediate stage have been established.

Let $\PRS(k)$ denote the projective Reed--Solomon code of length $q+1$
obtained by evaluating homogeneous binary forms of degree $k-1$ on
$\PP^1(\F_q)$.  If the redundancy is
\[
  r=q+1-k,
\]
a parity-check matrix has as columns the $\F_q$-points of the
degree-$(r-1)$ normal rational curve
\[
  \NRC_{r-1}=\{(1,t,\ldots,t^{r-1}):t\in\F_q\}
              \cup\{(0,\ldots,0,1)\}\subset\PP^{r-1}.
\]
Call a projective syndrome direction \emph{split-free} if it is not
contained in the span of $r-2$ distinct columns.  When
$\rho(\PRS(q+1-r))=r-1$, the split-free directions are precisely the
projective syndrome directions of deep-hole cosets.  The logical
separation used throughout the paper is
\[
\boxed{\quad
 \text{split-free classification}
 \;+\;\rho(\PRS(q+1-r))=r-1
 \quad\Longrightarrow\quad
 \text{deep-hole classification}.
\quad}
\]
The first term is proved geometrically; the covering-radius equality
is a separate coding-theoretic input and is never inferred from a
finite split-free census.

\noindent\emph{Two logical layers.}
The arbitrary-redundancy result separates an unconditional component theorem
from a conditional finite-field escape theorem.  Propositions~\ref{prop:reduced-terminal-carrier}--%
\ref{prop:recursive-component-selection} identify the reduced recursively
contained locus without a field-size bound or a lower-package hypothesis.
The theorem below then assumes the pointed lower packages and its displayed
bound on \(q\) to show that every split-free point lies in that locus.  The
fixed-level theorems discharge those extra inputs only through redundancy
ten.

\begin{theorem}[Recursive obstruction and fixed-level consequences]
\label{thm:main}
Let \(r\geq6\), and let
\[
 q\geq6r-15+\left\lfloor2\sqrt{6r-17}\right\rfloor .
\]
Every split-free redundancy-\(r\) syndrome belongs to
\[
 \cP_r\cup
 \PP\!\left\langle e_j:
 \binom{r-2}{j}=\binom{r-2}{j-1}=0
 \pmod{\operatorname{char}\F_q}\right\rangle .
 \autoeqtag{eq:main-carrier-containment}
\]
provided the explicit one-step lower packages of
Definition~\ref{def:lower-package}, with the stated deletion bounds, exist at
every intermediate redundancy.
Under the same package hypothesis, if
\(\operatorname{char}\F_q>r-1\), the second term is empty and the
deep-hole syndromes are exactly the persistent tangent and conjugate-secant
families, of total cardinality \(q(q+1)^2/2\).

At fixed levels, the classifications are exact at redundancies five and six
for every \(q\geq7\), at redundancy seven in the split-free sense for every
\(q\geq7\), and at redundancies eight, nine, and ten for
\(q\geq43,53,59\), respectively.  The redundancy-seven result promotes to
deep holes whenever the radius is six, in particular for \(q\geq11\).
\end{theorem}

Theorems~\ref{thm:r5}, \ref{thm:r6}, \ref{thm:r7}, \ref{thm:r8},
\ref{thm:r9}, and Corollary~\ref{cor:r10} give the fixed levels.  The
redundancy-ten conclusion uses Theorem~\ref{thm:m9-shallow}: the first fresh
higher Lucas carrier is entirely shallow, including the two-dimensional
quotient missed by the projective-additive section.

The correspondence between deep holes,
one-symbol MDS extensions, and uncovered points of a normal rational
curve was developed in
\cite[Prop.~1]{Kaipa2017} and
\cite[Lem.~II.4 and Thms.~I.4--I.6]{ZWK2020}; it belongs to the regime
in which the covering radius is \(r\), and every radius proved here is
\(r-1\), where D\"ur's equivalence instead forbids the extension.  See
Remark~\ref{rem:q8-no-extension}.
Zhang, Wan, and Kaipa classified projective Reed--Solomon deep-hole
directions through redundancy four and identified redundancy five as
the next case \cite{ZWK2020}.  Theorem~\ref{thm:r5} closes that case;
Theorem~\ref{thm:spine} gives the complete redundancy-six
classification and the radius-delimited redundancy-seven result.
Their lower tangent and conjugate-secant families remain inputs to the
terminal redundancy-five analysis.  Once that terminal classification is
known, the finite-depth theorem propagates it through marked contractions;
the prior lower-redundancy theory thus supplies the base data for the new
recursive mechanism.

Thus redundancy five is the first previously open level in the sequence,
and Theorem~\ref{thm:r5} classifies it over every prime power \(q\geq7\).

\begin{theorem}[Complete redundancy-five classification]\label{thm:r5}
Let $q\geq7$ be a prime power.  The covering radius of $\PRS(q-4)$ is
$4$.  Its projective deep-hole syndromes are exactly the following
disjoint families:
\begin{enumerate}[label=\textup{(\roman*)}]
\item the tangent family $\mathrm T$, of size $q(q+1)$;
\item the conjugate-secant family $\mathrm S$, of size
      $q(q^2-1)/2$;
\item if $\operatorname{char}\F_q\ne3$, the rational osculating-pair
      family $\mathrm O^+$ when $q\equiv2\pmod3$, of size $q(q+1)/2$,
      or the conjugate osculating-pair family $\mathrm O^-$ when
      $q\equiv1\pmod3$, of size $q(q-1)/2$;
\item in characteristic three, one $\PGL_2(q)$-fixed nucleus point and
      one wild Artin--Schreier orbit $\mathrm W$ of size $(q^2-1)/2$;
\item the sporadic $\PGL_2(q)$-orbits in Table~\ref{tab:r5sporadic},
      occurring exactly for
      $q\in\{7,8,9,11,13,17,19\}$.
\end{enumerate}
There are no other deep-hole syndromes.  The three corresponding
nonsporadic totals are
\[
\frac{q^2(q+3)}2,\qquad
\frac{q(q+1)(q+2)}2,\qquad
\frac{q^3+3q^2+q+1}{2},
\]
according as $q\equiv1\pmod3$ with characteristic not three,
$q\equiv2\pmod3$, or $\operatorname{char}\F_q=3$.
\end{theorem}

At \(q=8\) the classification gives \(1116\) deep-hole directions,
splitting as \(360\) nonsporadic and \(756\) sporadic ones.

The Hankel dictionary turns redundancy five into a pencil of cubics.
Cyclic descent gives the arithmetic families in
Theorem~\ref{thm:r5}; the \(S_3\) case is controlled by the
off-diagonal fiber square and a genus-one point bound.  At higher
redundancy, coherent polar contraction lowers the degree while retaining the
removed root as a forbidden marker.  The recursive carrier theorem identifies
the polar lines that remain trapped; the fixed-level arguments evaluate the
transverse packages and Lucas-carrier arithmetic.  Section~\ref{sec:overview}
gives the short proof route.

\begin{theorem}[Redundancy-six deep holes and redundancy-seven split-free syndromes]
\label{thm:spine}
The following statements hold.
\begin{enumerate}[label=\textup{(\alph*)}]
\item For every prime power $q\geq7$, the projective deep-hole
      directions of $\PRS(q-5)$ are exactly the persistent tangent and
      conjugate-secant families, the five finite-field exception rows of
      Theorem~\ref{thm:r6}, and, when $q=2^m$ with odd $m\geq5$, one
      characteristic-two nucleus orbit.
\item For every prime power $q\geq7$, Theorem~\ref{thm:r7} gives the
      complete split-free redundancy-seven classification.  For $q\geq13$
      only the persistent families and, when $q=2^m$ with odd $m$, one
      central nucleus point remain.  The result is a deep-hole classification
      for $q\geq11$.  At $q=8$ the known radius is seven; exact extraction
      leaves the nine-direction diagonal tangent orbit and the fixed central
      nucleus, so the classification is complete there as well. The radii at
      $q=7,9$ remain separate.
\end{enumerate}
\end{theorem}

Table~\ref{tab:classification-status} separates the split-free result,
radius gate, deep-hole conclusion, and remaining field gap at every level.
The R6 and R7 conclusions are all-field structural classifications: their
polar arguments exhaust every field outside a bounded range, and deterministic
certificates close the remaining prime powers.
